% Part:sets-functions-relations
% Chapter: size-of-sets
% Section: introduction

\documentclass[../../../include/open-logic-section]{subfiles}

\begin{document}

\olfileid{sfr}{siz}{int}

\olsection{پېژندنه}

کله چې ګېورګ کانتور په 1870مو کلونو کښې د سټونو نظريې ته وده ورکړه،
د هغه يو مقصد دا و چې د نامتناهي مجموعې مفکوره د منلو وړ کړي---
هغه څه چې د منځنيو پېړيو پوهانو به فعلي نامتناهي والے بللو۔ د
دې يوه مهمه برخه د بېلو سټونو د \emph{اندازې} په اړه د هغه
څېړنه وه۔ که $a$، $b$ او $c$ ټول بېل وي، نو په وجداني ډول
سټ $\{a, b, c\}$ له $\{a, b\}$ څخه \emph{لوے} دے۔ خو د
نامتناهي سټونو په اړه څه وايو؟ ايا ټول د يو بل هومره لوے دي؟
معلومېږي چې داسې نۀ ده۔

دلته لومړۍ مهمه مفکوره د شمېرنې ده۔ هر متناهي سټ د هغه د ټولو
\psOblique{!!{element}s} په ليکلو په لړ کښې راوړلے شو۔ د ځينو
نامتناهي سټونو ټول !!{element}s هم په لړ کښې راوړلے شو، که
اجازه ورکړو چې لړ پخپله نامتناهي وي۔ داسې سټونه !!{enumerable}
بلل کېږي۔ د کانتور حيرانوونکې پايله، چې د دې باب تر پايه به
ئې پوره وپېژنو، دا وه چې ځينې نامتناهي سټونه !!{enumerable} نۀ دي۔

\end{document}
